\documentclass{beamer}
\beamertemplatenavigationsymbolsempty
\usecolortheme{beaver}
\setbeamertemplate{blocks}[rounded=true, shadow=true]
\setbeamertemplate{footline}[page number]
%
\usepackage[utf8]{inputenc}
\usepackage[english,russian]{babel}
\usepackage{amssymb,amsfonts,amsmath,mathtext}
\usepackage{subfig}
\usepackage[all]{xy} % xy package for diagrams
\usepackage{array}
\usepackage{multicol}% many columns in slide
\usepackage{hyperref}% urls
\usepackage{hhline}%tables
\usepackage{natbib}
\usepackage{natbib}
\usepackage{doi}
\usepackage{mathtools}
% Your figures are here:
\graphicspath{ {fig/} {../fig/} }

%----------------------------------------------------------------------------------------------------------
\title{Дистилляция разнородных моделей в глубоком обучении}
\author[М.\,Ф. Набиев]{
    Набиев Мухаммадшариф Фуркатович \\
    Научный руководитель: к.ф.-м.н. О.\,Ю.~Бахтеев
}
\institute[Московский физико-технический институт]{
\small{
    Московский физико-технический институт \\
    Кафедра интеллектуальных систем ФПМИ МФТИ 
}}
\date{2025}
%----------------------------------------------------------------------------------------------------------
\begin{document}
%----------------------------------------------------------------------------------------------------------
\begin{frame}
\thispagestyle{empty}
\maketitle
\end{frame}
%---------------------------------------------------------------------------------------------------
\begin{frame}{Цель исследования}
\textbf{Проблема:} Методы дистилляции для \emph{разнородных} архитектур (например, RNN $\leftrightarrow$ CNN) обычно сопоставляют отдельные признаки, но слабо учитывают \textbf{распределение} скрытых представлений и их \textbf{поток} по слоям.
\newline 
\newline 
\textbf{Цель:} Предложить метод дистилляции разнородных моделей учитывающий распределение скрытых представлений.
\newline
\newline
\textbf{Решение:} Свести набор скрытых состояний к компактным матричным статистикам и выровнить их учитывая стоимость и вес переноса.
\end{frame}
%-------------------------------------------------------------------------------------------------- 
\begin{frame}{Постановка задачи}
\small
Пусть \(\mathfrak{D} = \{(\mathbf{x}_i, y_i)\}_{i=1}^n\)~-- набор изорбражений. Рассмотрим две модели \(R\) (RNN) и \(C\) (CNN). Их скрытые представления:
\[
    \mathbf{h}_1^R, \dots, \mathbf{h}_K^R \in \mathbb{R}^{D_1} \quad \text{и} \quad 
    \mathbf{A}_1^C, \dots, \mathbf{A}_N^C \in \mathbb{R}^{D_2}
\]
Для аггрегации скрытых переменных введем 
\[
    \mathbf{F}_k^R = \mathbf{u}_k\mathbf{u}_k^T \in \mathbb{R}^{d\times d}, \; \text{где} \; \mathbf{u}_k = \operatorname{aggr}(\mathbf{h}^R_1, \dots, \mathbf{h}^R_K) \in \mathbb{R}^{d}
\]
\[
    \mathbf{F}_j^C = \frac{1}{P_j} \mathbf{B}_j \mathbf{B}_j^T \in \mathbb{R}^{d\times d}, \; \text{где} \; \mathbf{B}_j = \operatorname{Flatten(}\operatorname{Conv}(\mathbf{A}_j^C)) \in \mathbb{R}^{d \times P_j}
\]
Имея набор аггрегированных матриц определеим матрицу стоимости \(\mathbf{C}\in \mathbf{R}^{K\times N}\), как \(C_{kj} = \| \mathbf{F}_k^R - \mathbf{F}_j^C \|_F^2\) и матрицу весов \(\boldsymbol{\Gamma} = \operatorname{Softmax(- \mathbf{C})}\).
Минимизируем следуюущую функцию потерь, 
\[
    \mathcal{L} = \mathcal{L} + \lambda \mathbb{E}_{\mathbf{x}} \sum_{k=1}^K \sum_{j=1}^N \Gamma_{kj}(\mathbf{x}) C_{kj}(\mathbf{x})
\]
В зависимости от выбора учителя (RNN или CNN) можно нормировать $\boldsymbol{\Gamma}$ по строкам или по столбцам.
\end{frame}
%----------------------------------------------------------------------------------------------------------
\begin{frame}{FGMKD: цель исследования}
\begin{block}{Проблема}
На практике при дистилляции есть компромисс: чем сильнее мы заставляем студента копировать учителей (\textbf{fidelity}), тем не всегда лучше \textbf{обобщение} (generalization).
\end{block}

\begin{block}{Цель}
Исследовать зависимость fidelity и generalization при дистилляции с ансамблем учителей.
\end{block}

\begin{block}{Предложенный подход}
Агрегировать учителей и обучать студента с коэффициентом $\lambda$ перед KL, чтобы получать фронт fidelity--generalization по $\lambda$.
\end{block}
\end{frame}
%----------------------------------------------------------------------------------------------------------
\begin{frame}{FGMKD: постановка задачи}
\small
Имеем данные \(\mathfrak{D}=\{(\mathbf{x}_j,y_j)\}_{j=1}^{n}\),
ансамбль учителей \(\{T_i(\cdot\mid \mathbf{x})\}_{i=1}^{M}\) и студент \(S_\theta(\cdot\mid \mathbf{x})\).
Центроид учителей:
\[
\hat{T}(\cdot\mid\mathbf{x})=\frac{1}{M}\sum_{i=1}^{M}T_i(\cdot\mid\mathbf{x}).
\]
Fidelity-метрики:
\[
F_i(\theta)=\mathbb{E}_{\mathbf{x}}\mathrm{KL}\!\left(T_i(\cdot\mid\mathbf{x})\ \|\ S_\theta(\cdot\mid\mathbf{x})\right),\quad
F_{\mathrm{avg}}(\theta)=\frac{1}{M}\sum_{i=1}^{M}F_i(\theta),
\]
\[
F_{\mathrm{cent}}(\theta)=\mathbb{E}_{\mathbf{x}}\mathrm{KL}\!\left(\hat{T}(\cdot\mid\mathbf{x})\ \|\ S_\theta(\cdot\mid\mathbf{x})\right).
\]
Diversity (разброс учителей) и декомпозиция:
\[
D_{\mathrm{div}}=\mathbb{E}_{\mathbf{x}}\frac{1}{M}\sum_{i=1}^{M}\mathrm{KL}\!\left(T_i(\cdot\mid\mathbf{x})\ \|\ \hat{T}(\cdot\mid\mathbf{x})\right),
\qquad
F_{\mathrm{avg}}(\theta)=D_{\mathrm{div}}+F_{\mathrm{cent}}(\theta).
\]
\end{frame}
%---------------------------------------------------------------------------------------------------
\begin{frame}{FGMKD: эксперименты}
\small
\begin{figure}[h!]
    \centering
  \begin{minipage}{0.9\textwidth}
    \includegraphics[width=\linewidth]{sergey_1_exp.jpg}
    
  \end{minipage}\hfill
\end{figure}
Обучение с центроидом дает прирост в качестве и в уменьшении \(\operatorname{KL}\).
\end{frame}
%----------------------------------------------------------------------------------------------------------
\begin{frame}{FGMKD: эксперименты}
\small
\begin{figure}[h!]
    \centering
  \begin{minipage}{0.9\textwidth}
    \includegraphics[width=\linewidth]{fideity_vs_gen_best-3.png}
  \end{minipage}\hfill
\end{figure}
Получили подтверждение оценки: 
\begin{equation}
\label{eq:risk_bound_main}
\bigl| R_{\mathrm{sup}}(S_\theta) - R_{\mathrm{sup}}(\hat{T}) \bigr|
\;\le\;
F_{\mathrm{cen}}(\theta)+C(\delta)\,\varepsilon \,\sqrt{F_{\mathrm{cen}}(\theta)},
\end{equation}
для подходящих констант $\detlta \in (0, 1)$ и $\varepsilon \ge 0$.
\end{frame}
%----------------------------------------------------------------------------------------------------------
\begin{frame}{Дальнейшие исследования}
Дистиллляций разнородных моделей:
\begin{itemize}
    \item Провести эксперименты на задаче классификации.
    \item Исследование влияний разных функций аггрегаций.
    \item Дистилляция в двух направлениях.
\end{itemize}
FGMKD:
\begin{itemize}
    \item Экспеименты на более сложных выборках.
    \item Взвешенные голосование учителей.
    \item Исследовать влияние плохих учителей.
    \item Учителя с ограниченным знанием (не знаю про существование остальных классов).
    \item Добавление шума в данных.
\end{itemize}
\end{frame}
%----------------------------------------------------------------------------------------------------------
% \begin{frame}{Источники}
%     \bibliographystyle{plain}
%     \bibliography{ib_in_model_selection}
% \end{frame}
%-----------------------------------------------------------------------------------------------------


\end{document} 
